\documentclass[10pt,a4paper]{amsart}
\usepackage[margin=19mm]{geometry}
\usepackage{amsmath,amssymb,mathtools}
\usepackage[scr=rsfs,cal=euler]{mathalfa}
\usepackage{xfrac}
\usepackage[unicode,hidelinks,bookmarks=false]{hyperref}
\newcommand{\ie}{i.e.\ }
\newcommand{\cf}{cf.\ }
\newcommand{\resp}{respectively\ }
\newcommand{\Subsection}{\subsection}
\providecommand{\nobreakdash}{\nobreak\hbox{-}}
\setlength{\emergencystretch}{2em}
\multlinegap0pt
\numberwithin{equation}{section}
\newcommand{\GG}{\mathcal{G}}
\newcommand{\FF}{\mathcal{F}}
\newcommand{\Tau}{\mathcal{T}}
\newcommand{\LL}{\mathcal{L}}
\newcommand{\clR}{\mathcal{R}}
\newcommand{\clG}{\mathcal{G}}
\newcommand{\MM}{\mathcal{M}}
\newcommand{\B}{\mathcal{B}}

\newcommand{\Z}{\mathbb{Z}}
\newcommand{\R}{\mathbb{R}}
\newcommand{\T}{\mathbb{T}}
\newcommand{\Prob}{\mathbb{P}}

\newcommand{\GL}{\mathrm{GL}}
\newcommand{\id}{\mathrm{id}}
\newcommand{\Diff}{\mathrm{Diff}}
\newcommand{\width}{\mathrm{width}}
\newcommand{\length}{\mathrm{length}}
\newcommand{\rspan}{\mathrm{span}}
\newcommand{\rst}{\mathrm{st}}
\newcommand{\bh}{\mathfrak{h}}
\newcommand{\bu}{\mathfrak{u}}
\newcommand{\bv}{\mathfrak{v}}


\newcommand*{\mk}{\mkern -1mu}
\newcommand*{\Mk}{\mkern -2mu}
\newcommand*{\mK}{\mkern 1mu}
\newcommand*{\MK}{\mkern 2mu}

%\hypersetup{urlcolor=purple, linkcolor=blue, citecolor=red}

\newcommand*{\romanenumi}{\renewcommand*{\theenumi}{\roman{enumi}}}
\newcommand*{\Romanenumi}{\renewcommand*{\theenumi}{\Roman{enumi}}}
\newcommand*{\alphenumi}{\renewcommand*{\theenumi}{\alph{enumi}}}
\newcommand*{\Alphenumi}{\renewcommand*{\theenumi}{\Alph{enumi}}}
\let\oldtilde\tilde
\renewcommand*{\tilde}[1]{\mathchoice{\widetilde{#1}}{\widetilde{#1}}{\oldtilde{#1}}{\oldtilde{#1}}}
\let\oldhat\hat
\renewcommand*{\hat}[1]{\mathchoice{\widehat{#1}}{\widehat{#1}}{\oldhat{#1}}{\oldhat{#1}}}
\renewcommand{\ie}{i.e.}
\newtheorem{theo}{Theorem}[section]
\newtheorem{defi}[theo]{Definition}
\newtheorem{lemm}[theo]{Lemma}
\newtheorem{prop}[theo]{Proposition}
\newtheorem{rema}[theo]{Remark}
\newtheorem{coro}[theo]{Corollary}
\title[Uniform width convergence at irrational directions]{Deforming a Finsler metric on the two-torus to a flat Finsler metric: original Theorem 5.1}
\author{Elie Nakhlé and Stéphane Sabourau}
\date{Source: Annales Henri Lebesgue 7 (2024), 1131-1174; DOI 10.5802/ahl.217}
\begin{document}
\maketitle
\noindent\textit{Editorial scope.} The counted unit is original Theorem 5.1 with the authors complete proof. The definition, remark and the two lemmas of original Section 5 that the proof invokes are retained as uncounted context. The rest of the article is omitted. One illustration float inside the proof is omitted by an explicit fidelity receipt and its reference is delivered as a hyperlink to the published Figure 5.1; no artwork is redistributed. The source's own commented-out alternative proofs are not delivered, exactly as in the published article. Original theorem, lemma, remark and definition numbers are retained. Only the AMS wrapper supplies editorial formatting.
\setcounter{section}{4}
\section{Convergence of the width at irrational directions} \label{sec:uniform}


We show that the width of a geodesic in the universal cover of a Finsler two-torus without conjugate points converges to zero under the curve shortening flow for geodesics whose asymptotic direction lies in a small enough neighborhood of a fixed irrational direction. We deduce that the straight lines obtained at the limit vary continuously at irrational directions.



\begin{theo}\label{theo:unifwidth}
Let $M=(\T^2,F)$ be a Finsler two-torus without conjugate points. Denote by~$\bar{M}=(\R^2,\bar{F})$ its universal Finsler cover. Let $\theta$ be an irrational direction. For every $\varepsilon>0$, there exists $t_0>0$ such that for every geodesic~$\gamma$ of~$\bar{M}$ with asymptotic direction close enough to~$\theta$ and every~$t \geq t_0$, the width of~$\gamma_t$ satisfies $\width(\gamma_t) < \varepsilon$.
\end{theo}

\begin{rema}
The point of Theorem~\ref{theo:unifwidth} is that $t_0$ does not depend on~$\gamma$ as long as its asymptotic direction remains close enough to the irrational direction~$\theta$. Of course, Theorem~\ref{theo:unifwidth} also applies when the asymptotic direction of the geodesic~$\gamma$ is irrational, which gives an analogue of Lemma~\href{https://ahl.centre-mersenne.org/item/AHL_2024__7__1131_0/}{4.9} in the irrational case and allows us to derive Theorem~\href{https://ahl.centre-mersenne.org/item/AHL_2024__7__1131_0/}{4.4} in this case too. We will need this stronger version of Theorem~\href{https://ahl.centre-mersenne.org/item/AHL_2024__7__1131_0/}{4.4} to establish the continuity of the Euclidean curve shortening flow at the limit for Finsler geodesics with irrational asymptotic directions.
\end{rema}

We decompose the proof of the proposition into several lemmas.

\begin{defi}
The \emph{vertical distance} between two subsets $A, B \subset \R^{2}$ is defined as
\[
\sup_{V} \inf\{d(x,y) \mid x\in A\cap V, y\in B\cap V\}
\]
where $V$ runs over all vertical lines of~$\R^2$ intersecting~$A$ and~$B$, and $d$ represents the Euclidean distance in $\R^{2}$.
\end{defi}

The following lemma is a quantitative version of the fact that an irrational geodesic in the square flat torus is dense.

\begin{lemm}\label{lem:dioph}
Let $\theta$ be an irrational direction. For every $\varepsilon \in (0,1)$, there exists $L=L(\theta,\varepsilon)>0$ such that every segment~$c$ of~$\R^2$ of direction close enough to~$\theta$ which projects onto an interval of the $x$-axis of length at least~$L$ passes below a point~$x_+$ of~$\Z^2$ and above a point~$x_-$ of~$\Z^2$, both at vertical distance less than~$\varepsilon$ from the segment~$c$.
\end{lemm}


\begin{proof}
Let $\varepsilon \in (0,1)$. Fix a point~$\bar{x} \in \Z^2$. Let $\bar{x}_+$ and~$\bar{x}_-$ be the points of~$\R^2$ at vertical distance~$\varepsilon$ from~$\bar{x}$, with $\bar{x}_+$ above~$\bar{x}$ and $\bar{x}_-$ below~$x$. Let~$\bar{m}_+$ and~$\bar{m}_-$ be the midpoints of~$[\bar{x},\bar{x}_+]$ and~$[\bar{x},\bar{x}_-]$. Denote by~$\pi(\bar{m}_+)$ and~$\pi(\bar{m}_-)$ the projections of~$\bar{m}_+$ and~$\bar{m}_-$ to~$\T^2$, where $\pi:\R^2 \to \T^2$ is the covering projection. Suppose that the segment~$c$ is directed along the irrational direction~$\theta$. If the segment~$c$ is long enough (depending on~$\theta$ and~$\varepsilon$), its projection to~$\T^2$ passes at (vertical) distance less than~$\frac{\varepsilon}{2}$ from any point, and in particular from~$\pi(\bar{m}_+)$ and~$\pi(\bar{m}_-)$. Thus, the segment~$c$ passes at vertical distance less than~$\frac{\varepsilon}{2}$ from some $\Z^2$-translates~$m_+$ and~$m_-$ of~$\bar{m}_+$ and~$\bar{m}_-$ (\ie, some lifts of~$\pi(\bar{m}_+)$ and~$\pi(\bar{m}_-)$). Define~$x_+$ and~$x_-$ as the corresponding $\Z^2$-translates of~$\bar{x}_+$ and~$\bar{x}_-$. By construction, the segment~$c$ passes below~$x_+$ and above~$x_-$ at vertical distance less than~$\varepsilon$ from them. The same holds for a segment~$c$ of direction close enough to $\theta$.
\end{proof}

Denote by~$(\vec{\imath},\vec{\jmath})$ the canonical basis of~$\R^2$. The following lemma is a quantitative version of the fact that the curve shortening flow straighten curves.

%\forget



\begin{lemm}\label{lem:L'}
Assume that the horizontal and vertical curves of~$M$ are geodesic. For every $\varepsilon \in (0,1)$ and every~$L>0$, there exists $t_0=t_0(\varepsilon,L)>0$ such that for every $t \geq t_0$ and every geodesic~$\gamma$ of~$\bar{M}$ whose asymptotic direction forms an angle lying between~$-\frac{\pi}{4}$ and~$\frac{\pi}{4}$ with the horizontal direction, the following holds: every arc of~$\gamma_t$ over an interval~$I$ of length~$L$ of the $x$-axis is at vertical distance at most~$\varepsilon$ from a segment of~$\R^2$ over the same interval~$I$.
\end{lemm}

\begin{proof}
By Theorem~\href{https://ahl.centre-mersenne.org/item/AHL_2024__7__1131_0/}{2.7}, the geodesic~$\gamma$ lies in a strip~$S$ of width~$w$ with the same asymptotic direction, where $w$ does not depend on~$\gamma$ (only on the Finsler metric on~$M$). Since the vertical curves of~$\bar{M}$ are geodesic, the tangent vector~$\gamma'$ is uniformly bounded away from a vertical direction, that is, $|\langle \gamma',\vec{\jmath} \rangle| \leq C \, \lVert \gamma' \rVert$ for some constant $C \in (0,1)$ not depending on the geodesic~$\gamma$. Otherwise, the geodesic~$\gamma$ would be close to a vertical geodesic and would leave the strip~$S$ whose vertical width is bounded.

By Lemma~\href{https://ahl.centre-mersenne.org/item/AHL_2024__7__1131_0/}{4.7}, there exists~$c_0>0$ not depending on~$\gamma$ such that the curvature~$\kappa(\gamma_t)$ of~$\gamma_t$ satisfies
\[
|\kappa(\gamma_t)| \leq \frac{c_0}{\sqrt{t}}.
\]
Thus, for $t$ large enough, the curve~$\gamma_t$ is almost straight. Hence the desired result.
\end{proof}

We can now proceed to the proof of Theorem~\ref{theo:unifwidth}.


\begin{proof}[Proof of Theorem~\ref{theo:unifwidth}]
By Proposition~\href{https://ahl.centre-mersenne.org/item/AHL_2024__7__1131_0/}{3.1}, we can assume that the horizontal and vertical curves of~$\T^2$ are geodesics. Switching the roles of the $x$- and~$y$-axis, and the orientation of~$\gamma$ if necessary, we can further assume that the angle between the asymptotic direction of~$\gamma$ and the horizontal vector~$\vec{\imath}$ lies between~$-\frac{\pi}{4}$ and~$\frac{\pi}{4}$. Recall that the curve~$\gamma_t$ lies in a (closed) Euclidean strip~$S_t=S(\gamma_t)$ of~$\R^2$ bounded by two straight lines~$\Delta_+=\Delta_+(\gamma_t)$ and~$\Delta_-=\Delta_-(\gamma_t)$ with the same asymptotic direction as~$\gamma$, with $\Delta_+$ above~$\Delta_-$ in~$\R^2$; see Theorem~\href{https://ahl.centre-mersenne.org/item/AHL_2024__7__1131_0/}{2.7}.

Let us show that if the asymptotic direction of~$\gamma$ is close enough to~$\theta$ (more precisely, if the asymptotic direction of~$\gamma$ is in the neighborhood of~$\theta$ given by Lemma~\ref{lem:dioph}), then the width of~$S_t$ is less than~$14\varepsilon$ for $t \geq t_0$. Take a point~$p_+=p_{+}(\gamma_t)$ of~$\gamma_t$ at vertical distance at most~$\varepsilon$ from~$\Delta_+$ (below~$\Delta_+$) and denote by~$\bar{p}_+$ its projection to the $x$-axis. Let $I$ be the interval of the $x$-axis centered at~$\bar{p}_+$ of length~$L$, where $L$ is given in Lemma~\ref{lem:dioph}. Consider the arc~$\alpha_+$ of~$\gamma_t$ over~$I$. By Lemma~\ref{lem:L'}, this arc~$\alpha_+$ is at vertical distance at most~$\varepsilon$ from a segment~$c$ of~$\R^2$. Thus, the segment~$c$ lies below the line~$\Delta_+ +\varepsilon \vec{\jmath}$ above~$\Delta_+$ at vertical distance~$\varepsilon$ from~$\Delta_+$ and passes through a point at vertical distance at most~$\varepsilon$ from~$p_+$. Therefore, the segment~$c$ lies at vertical distance at most~$6\varepsilon$ from~$\Delta_++\varepsilon \vec{\jmath}$. We deduce that the arc~$\alpha_+$ of~$\gamma_t$ lies below~$\Delta_+$ at vertical distance at most~$6\varepsilon$ from~$\Delta_+$. See~Figure~\href{https://ahl.centre-mersenne.org/item/AHL_2024__7__1131_0/}{5.1}. 



Apply Lemma~\ref{lem:dioph} to the segment lying in~$\Delta_+ - 6 \varepsilon \vec{\jmath}$ above the interval~$I$ of length~$L$, assuming the asymptotic direction of~$\Delta_+$ is close enough to the irrational direction~$\theta$. This yields a point~$x_-$ of~$\Z^2$ below~$\Delta_+-6 \varepsilon \vec{\jmath}$ at vertical distance at most~$\varepsilon$ from it and so at vertical distance at most~$7\varepsilon$ from~$\Delta_+$. It follows that the geodesic arc~$\alpha_+$ of~$\gamma_t$ passes above the point~$x_-$ of~$\Z^2$ at vertical distance at most~$7\varepsilon$ from~$\Delta_+$. Similarly, the curve~$\gamma_t$ passes below a point~$x_+$ of~$\Z^2$ at vertical distance at most~$7\varepsilon$ from~$\Delta_-$.

Let $T$ be the translation of~$\R^2$ by the integral vector~$x_--x_+ \in \Z^2$. Note that $T$ is an isometry both for the Euclidean metric and the Finsler metric. By construction, the point~$x_-$ lies below~$\gamma_t$ and above $T(\gamma_t)$, and is at vertical distance at most~$7\varepsilon$ from~$\Delta_+$ and~$T(\Delta_-)$. The geodesics~$\gamma$ and~$T(\gamma)$ have the same asymptotic direction and do not intersect. The same holds for their images~$\gamma_t$ and~$T(\gamma_t)$ under the curve shortening flow; see Section~\href{https://ahl.centre-mersenne.org/item/AHL_2024__7__1131_0/}{4}. In addition, $T(\gamma_{t})$ is proved to be below $\gamma_{t}$ and $T(\Delta_{-})$ is below $\Delta_{-}$, hence $T(\Delta_{-})$ bounds $\gamma_{t}$ from below. It follows that the curve~$\gamma_t$ lies in the strip of width at most~$14 \varepsilon$ bounded by~$\Delta_+$ on top and by~$T(\Delta_-)$ at the bottom. Thus, $\width(\gamma_t) \leq 14 \varepsilon$.
\end{proof}

\end{document}
